% TOP_PROBLEM: {"id":30007521,"problem_number":"LOCAL-30007521","title":"Structural change in cyclical models","category_id":21,"category":{"id":21,"name":"economics","display_name":"Economics","description":"Open economic theory statements, published research agendas, and proposed scoped specifications; question provenance and historical priority are qualified per record.","slug":"economics","order_index":21,"created_at":"2026-10-08T00:00:00Z"},"status":"research_question","external_url":null,"legacy_ids":[],"published":false,"statement":"Identify which climate, digital-capital, and demographic channels improve policy-response predictions in a declared nested macroeconomic model.","economics":{"original_id":10,"field":"Business cycles and macro dynamics","kind":"empirical","formalization_basis":"adapted_research_question","source_status":"research_agenda","priority_status":"not_established","priority_note":"Original priority is not established. No dated primary statement of this precise problem was verified.","earliest_verified_reference":{"authors":["European Central Bank"],"title":"Research agenda: Forecasting, macroeconomic and financial analysis","year":null,"url":"https://www.ecb.europa.eu/pub/economic-research/research_agenda/forecasting/html/index.en.html","doi":null,"locator":"“Model development,” selected areas of focus","evidence_summary":"The agenda explicitly calls for models incorporating climate change, digitalisation, globalisation, demographic changes, and dynamic multicountry analysis."},"current_reference":{"authors":["European Central Bank"],"title":"Research agenda: Forecasting, macroeconomic and financial analysis","year":null,"url":"https://www.ecb.europa.eu/pub/economic-research/research_agenda/forecasting/html/index.en.html","doi":null,"locator":"“Model development,” selected areas of focus","evidence_summary":"The agenda explicitly calls for models incorporating climate change, digitalisation, globalisation, demographic changes, and dynamic multicountry analysis."},"formalization_note":"The agenda explicitly calls for these model extensions. The intervention and nested-model comparison are newly scoped; no canonical equation is attributed to the ECB."},"statusQualification":"Published question or research agenda verified at the cited locator. The mathematical specification is adapted for this catalogue and includes additional assumptions; source publication does not establish first-ever priority or certify this exact model remains unsolved. The agenda explicitly calls for these model extensions. The intervention and nested-model comparison are newly scoped; no canonical equation is attributed to the ECB.","statusReviewedAt":"2026-10-08"}
% FIRST_POSTED: 2026-10-08
% ENTRY_KIND: statement_only
% !TeX program = lualatex
\documentclass[11pt]{article}
\usepackage{fontspec}
\usepackage[a4paper,margin=25mm]{geometry}
\usepackage{amsmath,amssymb,mathtools}
\usepackage{xurl}
\usepackage[hidelinks]{hyperref}
\newcommand{\sref}[2]{\hyperlink{source:#1:#2}{[S#2]}}
\setlength{\emergencystretch}{3em}
\begin{document}
% problemId:30007521
\section*{Structural change in cyclical models}\label{30007521}
\subsection{Definitions and mathematical statement}
Fix a baseline macroeconomic model with equilibrium state \(s_t\), policy rule \(a_t\), and observable vector \(z_t\). Add measurable slow states \(v_t=(v_t^C,v_t^D,v_t^N)\), denoting climate exposure, digital-capital penetration, and demographic composition. Their definitions, units, data vintages, and transition laws must be specified. Let \(M_0\) hold those states fixed at a benchmark and let \(M_1\) allow them to affect production, demand, labor supply, and financial constraints through an identified finite-dimensional parameter \(\theta\).
For a common identified monetary or sectoral intervention \(a\), define
\[
\Delta(h,v)=E_{M_1}[z_{t+h}\mid\operatorname{do}(a_t=a),v_t=v]
-E_{M_0}[z_{t+h}\mid\operatorname{do}(a_t=a)].
\]
Determine which channels and interactions among the three slow states are required to predict \(\Delta(h,v)\) over a declared horizon and range of environments. Identify causal variation separately from concurrent trends, and avoid treating every slow movement as an independent shock. Compare nested restrictions using held-out economies, policy episodes, and forecast distributions; estimate the uncertainty attributable to each added channel. A successful result is a tractable, empirically supported extension with improved policy-response predictions under specified structural transitions. It is not a demand for one universal model covering all possible climate, technological, and demographic futures.

\par\textbf{Formulation and scope.} The agenda explicitly calls for these model extensions. The intervention and nested-model comparison are newly scoped; no canonical equation is attributed to the ECB.

\par\textbf{Open-question provenance.} An explicit question or research agenda was verified at the source locator. This does not by itself certify that every later version remains unresolved \sref{30007521}{1}.

\par\textbf{Historical priority.} Original priority is not established. No dated primary statement of this precise problem was verified.

\subsection{Short English statement}
Identify which climate, digital-capital, and demographic channels improve policy-response predictions in a declared nested macroeconomic model.

\subsection{Sources}
\begin{itemize}\raggedright
\item[\textnormal{[S1]}]\hypertarget{source:30007521:1}{} European Central Bank. Year not verified. Research agenda: Forecasting, macroeconomic and financial analysis. “Model development,” selected areas of focus.
\url{https://www.ecb.europa.eu/pub/economic-research/research_agenda/forecasting/html/index.en.html}.
{\small\textit{Earliest verified question/agenda reference; historical priority qualified below. The agenda explicitly calls for models incorporating climate change, digitalisation, globalisation, demographic changes, and dynamic multicountry analysis.}}
\end{itemize}
\end{document}
